\chapter{Step 4 — Deciding on Controls and Actions}
\label{chap:Step 4 — Deciding on Controls and Actions}

%---------------------------- SECTION ----------------------------
\section{From Understanding to Doing}
\paragraph{The first three steps of the risk assessment process have been, in essence, about understanding — understanding what risks exist, who they affect, and how significant they are. That understanding is valuable, but it is not the end point. A risk assessment that produces a thorough, well-evaluated picture of the risk landscape and then stops is, at best, an academic exercise. The value of everything that has come before is realised in this step — where understanding is converted into decisions, and decisions are converted into action.}
\paragraph{Step 4 is where the organisation decides what it is going to do about the risks it has identified and evaluated. It is where controls are selected, designed, and implemented — or where a conscious, documented decision is made that a risk is acceptable and no further action is required. It is where the risk assessment connects most directly to the day-to-day reality of how the organisation operates, and where the compliance professional's role shifts from analyst to adviser, from observer to participant.}
\paragraph{It is also, in the experience of most practitioners, where risk assessment most frequently falls short. Organisations that conduct thorough identification and careful evaluation sometimes struggle to translate that work into effective, proportionate, and genuinely implemented controls. The gap between knowing what the risks are and actually managing them well is where a great many regulatory failures, operational disasters, and reputational crises find their origin.}
\paragraph{This chapter is about closing that gap.}

%---------------------------- SECTION ----------------------------
\section{The Four Fundamental Options}
\paragraph{When faced with an evaluated risk, an organisation has four fundamental options. These are not mutually exclusive — in practice, most risk responses involve some combination of them — but understanding them as distinct options is the starting point for sound decision-making.}

\subsection{Option 1 — Avoid the Risk}
\paragraph{\textbf{Risk avoidance} means eliminating the risk entirely by choosing not to engage in the activity, process, or relationship that gives rise to it. If a proposed new product would create an unacceptable regulatory exposure, the organisation could choose not to launch it. If a particular supplier relationship carries conduct or reputational risks that exceed the organisation's appetite, the organisation could choose not to enter or to exit that relationship.}
\paragraph{Avoidance is the most complete form of risk response — it eliminates the exposure rather than merely managing it. It is also, however, the most commercially significant — because avoiding a risk frequently means forgoing the opportunity or benefit associated with the activity that creates it. For this reason, avoidance tends to be most appropriate where a risk is genuinely unacceptable — where it exceeds the organisation's risk appetite by a significant margin and where no other response can bring it within acceptable bounds}
\paragraph{It is worth noting that avoidance is not always possible or desirable even for very high risks. An organisation that responds to every high-rated risk by avoiding the underlying activity would quickly find itself unable to operate at all. The question is not simply \textit{can we avoid this risk? but is avoidance the most proportionate and appropriate response given the nature of the risk and the value of the activity that creates it?}}

\subsection{Option 2 — Reduce the Risk}
\paragraph{\textbf{Risk reduction} — sometimes called risk mitigation — means taking action to reduce either the likelihood of the risk materialising, the severity of its consequences if it does, or both. This is the most commonly deployed risk response option, and it encompasses the widest range of possible actions.}
\paragraph{Controls designed to reduce likelihood aim to prevent the risk from occurring — by strengthening processes, improving systems, enhancing training, increasing oversight, or addressing the underlying causes of the risk. Controls designed to reduce consequence aim to limit the harm if the risk does occur — through early detection, rapid response, business continuity arrangements, or the existence of financial buffers.}
\paragraph{For compliance professionals, the distinction between prevention and mitigation controls is practically important. A control that prevents a regulatory breach from occurring is preferable to one that simply detects it after the fact — but detection and response controls are a necessary complement to prevention, because no prevention control is perfect. A layered approach — combining controls that reduce likelihood with controls that limit consequence — is generally more robust than relying on any single measure.}

\subsection{Option 3 — Transfer the Risk}
\paragraph{\textbf{Risk transfer} means shifting some or all of the financial consequences of a risk to a third party. The most familiar mechanism for risk transfer is \textbf{insurance} — paying a premium to an insurer in exchange for protection against the financial consequences of defined risk events. Other mechanisms include contractual risk allocation — the use of indemnities, warranties, and liability caps in contracts to allocate risk between parties — and financial hedging instruments that transfer market risk.}
\paragraph{Risk transfer is an important and legitimate part of a comprehensive risk response strategy, but it comes with important caveats that compliance professionals should understand.}
\paragraph{\textbf{Transfer does not eliminate the risk}. Transferring the financial consequences of a risk to an insurer does not prevent the risk from materialising, does not protect the organisation from regulatory consequences, and does not prevent reputational damage. An organisation that experiences a significant data breach does not avoid regulatory investigation simply because it has cyber insurance — the ICO's assessment of the organisation's compliance will be entirely unaffected by the existence of an insurance policy.}
\paragraph{\textbf{Transfer may not be complete}. Insurance policies contain exclusions, sub-limits, and conditions that may mean the coverage available is less than assumed. Contractual risk allocation may be challenged or may prove unenforceable. Due diligence on the scope and reliability of transfer mechanisms is an important part of any risk response that relies significantly on transfer.}
\paragraph{\textbf{Some risks cannot be transferred}. Regulatory liability, criminal liability, and the direct consequences of regulatory breach — including loss of licence, formal censure, and personal liability for senior individuals — cannot be transferred to a third party. No insurance policy covers the reputational consequences of a major compliance failure. Some risks can only be managed, not transferred.}

\subsection{Option 4 — Accept the Risk}
\paragraph{\textbf{Risk acceptance} means making a conscious, documented decision that a risk is within the organisation's risk appetite and that no further action is required beyond ongoing monitoring. It is not the same as ignoring a risk — it is a deliberate, informed judgement that the cost and effort of further risk reduction is not proportionate to the benefit.}
\paragraph{Risk acceptance is entirely appropriate for low-rated risks that fall comfortably within the organisation's risk appetite. Indeed, an organisation that attempts to control every identified risk regardless of its significance is likely to waste significant resources and to dilute the focus of its risk management effort on the risks that actually matter.}
\paragraph{The key discipline around risk acceptance is that it must be \textbf{conscious and documented}. A risk that is accepted without acknowledgement — where the decision not to act is made by default rather than by deliberate judgement — is not managed acceptance. It is a gap. The difference between the two matters enormously if the risk subsequently materialises and the organisation's decision-making is scrutinised by a regulator, an auditor, or a court.}
\paragraph{Risk acceptance above a defined threshold should require the explicit sign-off of an appropriate level of management or governance — ensuring that the decision to live with a significant risk is made with full awareness and appropriate authority, not simply inherited as a default position.}

%---------------------------- SECTION ----------------------------
\section{The Hierarchy of Controls}
\paragraph{Choosing between and within these four options benefits from a structured framework — and the most widely recognised and practically useful framework for this purpose is the \textbf{Hierarchy of Controls}.}
\paragraph{Originally developed in the context of occupational health and safety — where it remains a central concept in regulatory frameworks including UK health and safety law — the Hierarchy of Controls has been adapted and applied across a wide range of risk management contexts. Its essential logic is simple and powerful: some types of control are inherently more reliable and more effective than others, and organisations should prefer higher-order controls where they are practicable.}
\paragraph{The hierarchy, from most to least effective, runs broadly as follows:}

\subsection{Level 1 — Elimination}
\paragraph{Eliminating the hazard entirely — removing the source of risk rather than managing exposure to it. This is the risk avoidance option described above, translated into control terms. It is the most effective form of control because it means the risk no longer exists.}
\paragraph{\textit{Example: Ceasing to collect a category of personal data that is not necessary for any business purpose eliminates the data protection risk associated with holding and processing that data.}}

\subsection{Level 2 — Substitution}
\paragraph{Replacing the hazardous activity, process, or element with a less risky alternative. The risk is not eliminated entirely, but its fundamental nature is changed in a way that reduces the level of exposure.}
\paragraph{\textit{Example: Replacing a manual, paper-based regulatory reporting process — with its inherent risks of human error and lost documentation — with an automated, system-based process that reduces the opportunity for error and creates an automatic audit trail.}}

\subsection{Level 3 — Engineering Controls}
\paragraph{Controls that are built into the design of a process, system, or environment — making the risky outcome harder to achieve through the structure of the activity itself, rather than relying on the behaviour of individuals.}
\paragraph{In a compliance context, engineering controls translate broadly into \textbf{system controls and process design} — controls that are embedded in technology or process architecture rather than dependent on people choosing to follow a procedure. System-enforced segregation of duties, automated compliance checks, mandatory approval workflows, and hard-coded system limits are all examples of engineering-type controls in a compliance setting.}
\paragraph{These controls are generally more reliable than those that depend on individual behaviour, because they do not require anyone to remember to do something or to choose to follow a rule. They work by making the compliant outcome the default or the only possible outcome.}

\subsection{Level 4 — Administrative Controls}
\paragraph{Controls that operate through policies, procedures, training, and oversight — relying on people to behave in accordance with defined rules and standards. Administrative controls are the most commonly deployed category of control in compliance and legal risk management — the policies, procedures, training programmes, and management oversight arrangements that constitute the bulk of most organisations' compliance frameworks.}
\paragraph{Administrative controls are important and necessary, but they are inherently less reliable than higher-order controls because they depend on human behaviour. People forget, misunderstand, cut corners, and sometimes deliberately circumvent procedures. This does not mean administrative controls are ineffective — properly designed and consistently reinforced, they can be highly effective — but it does mean that they should not be relied upon as the only layer of protection for significant risks, and that their effectiveness needs to be actively monitored and tested.}

\subsection{Level 5 — Detective Controls}
\paragraph{Controls that do not prevent a risk from materialising but are designed to detect it quickly when it does — enabling early intervention and limiting the consequences. Monitoring systems, exception reports, audit programmes, whistleblowing mechanisms, and regulatory reporting requirements all function, at least in part, as detective controls.}
\paragraph{Detective controls are an essential complement to preventive controls — because no preventive control is perfect, and early detection of a problem is almost always better than late detection. The value of a detective control depends critically on two things: the speed with which it detects a problem, and the quality of the response it triggers.}

\subsection{Level 6 — Protective Controls}
\paragraph{Controls that do not prevent or detect a risk but are designed to limit the consequences when it materialises — minimising the harm and enabling recovery. Business continuity arrangements, crisis management plans, insurance, financial reserves, and incident response protocols are all examples of protective controls.}
\paragraph{For compliance professionals, protective controls are particularly important in relation to regulatory risk. When a compliance failure occurs — as eventually it will in most organisations, however well-managed — the organisation's ability to respond quickly, effectively, and transparently can make a very significant difference to the regulatory outcome. A firm that detects a breach, self-reports promptly, takes immediate remedial action, and cooperates fully with the regulator is likely to be treated very differently from one that is slow to recognise a problem, fails to report it, and is perceived to be minimising or concealing it.}

%---------------------------- SECTION ----------------------------
\section{Selecting Proportionate Controls — The Cost-Benefit Dimension}
\paragraph{The Hierarchy of Controls tells us that some controls are inherently more effective than others. But choosing which controls to implement is not only a question of effectiveness — it is also a question of proportionality. The effort, cost, and disruption involved in implementing a control needs to be proportionate to the significance of the risk it is designed to manage.}
\paragraph{This proportionality principle appears explicitly in many regulatory frameworks. UK health and safety law, for example, requires risks to be reduced \textit{so far as is reasonably practicable} — a formulation that explicitly acknowledges a cost-benefit dimension. The GDPR's requirement to implement \textit{appropriate} technical and organisational measures similarly implies a proportionality assessment.}
\paragraph{In practice, proportionality assessment involves weighing several factors:}
\paragraph{\textbf{The significance of the risk} — a high-residual-risk finding warrants more substantial and more costly control investment than a low-risk one. The starting point for any proportionality assessment is the risk evaluation produced in Step 3.}
\paragraph{\textbf{The cost and effort of the control} — including not just the direct financial cost but the operational burden, the management time, and the potential disruption to business activities. Some controls are cheap and simple; others are expensive and complex. The cost of implementation needs to be weighed honestly against the risk reduction benefit.}
\paragraph{\textbf{The effectiveness of available options} — a relatively inexpensive control that reliably reduces a risk by 80\% may be better value than an expensive control that reduces the same risk by 90\%. Understanding the marginal benefit of additional control investment is part of making good control decisions.}
\paragraph{The residual risk if the control is not implemented — the relevant comparison is not between the cost of the control and the abstract level of the risk, but between the cost of the control and the expected cost of the harm that would result if the risk materialised without the control in place. Where the expected harm significantly exceeds the cost of the control, the proportionality case for implementing the control is strong.}
\paragraph{\textbf{Regulatory expectations} — in some areas, regulatory guidance or enforcement practice effectively sets a floor for control investment that goes beyond what a pure cost-benefit analysis might suggest. Where regulatory expectations are clear, they need to be reflected in control decisions regardless of the outcome of an internal proportionality assessment.}

%---------------------------- SECTION ----------------------------
\section{Designing Controls That Actually Work}
\paragraph{Selecting controls is one thing. Designing them so that they actually work in practice is another — and the gap between the two is where a significant proportion of compliance failures originate.}
\paragraph{Several principles are worth keeping in mind when designing controls.}

\subsection{Controls Must Be Specific and Actionable}
\paragraph{A control described as \textit{"management oversight of the process"} is unlikely to be effective — because it does not specify who is responsible, what they are supposed to do, how often they are supposed to do it, or what they should do if they identify a problem. Effective controls are specific: they describe a defined action, performed by a named role, at a defined frequency, with a defined output or outcome.}
\paragraph{The test of whether a control is sufficiently specific is whether someone reading the control description for the first time would know exactly what they need to do to operate it. If the answer is no — if the description leaves significant room for interpretation — the control needs further definition.}

\subsection{Controls Must Be Owned}
\paragraph{Every control needs a named owner — an individual who is responsible for ensuring that it is implemented, operated as designed, and maintained over time. Controls without clear ownership tend to be inconsistently applied, poorly maintained, and the first things to be deprioritised when capacity is under pressure.}
\paragraph{The owner of a control does not need to be the person who operates it day-to-day. A senior manager may own a control that is operated by their team — but the ownership means they are accountable for its effectiveness and will be held to account if it fails.}

\subsection{Controls Must Be Tested}
\paragraph{Documented controls and effective controls are not the same thing. A policy that nobody reads, a training programme that nobody retains, and a system control that has been misconfigured are all documented controls that provide no actual protection. Testing — through audit, monitoring, sampling, and periodic assessment — is the mechanism by which the organisation assures itself that its controls are working as intended.}
\paragraph{For compliance professionals, the testing of key controls is a core activity. It is also, increasingly, something that regulators specifically expect to see evidence of. A firm that can demonstrate not just that controls exist but that they have been tested, that they are operating effectively, and that identified weaknesses have been addressed is in a much stronger position than one that can only produce documentation.}

\subsection{Controls Must Be Resilient}
\paragraph{Controls that depend on a single person, a single system, or a single process are vulnerable to single points of failure. Where a control is critical — where its failure would allow a significant risk to materialise — building in redundancy and resilience is important. This might mean having a backup process for a critical system control, cross-training staff so that a key control can be operated in someone's absence, or designing escalation paths that do not depend on a single individual's availability or judgement.}

\subsection{Controls Must Be Proportionate to the Risk}
\paragraph{Returning to the proportionality principle — controls that are disproportionately burdensome relative to the risk they manage create their own problems. Over-controlled processes generate unnecessary friction, consume resources that could be better deployed elsewhere, and — perhaps most insidiously — breed a culture of compliance theatre, where people go through the motions of control activities without genuine engagement. Well-designed controls are efficient as well as effective.}

%---------------------------- SECTION ----------------------------
\section{The Control Framework — Building Coherence}
\paragraph{Individual controls do not exist in isolation. They sit within a broader \textbf{control framework} — the totality of the policies, procedures, systems, oversight arrangements, and cultural factors that together determine how well the organisation manages its risks.}
\paragraph{For compliance professionals, one of the most important tasks in the control design step is ensuring that the overall control framework is \textbf{coherent} — that individual controls fit together logically, that they do not contradict each other, that there are no significant gaps between them, and that they collectively address the organisation's risk profile in a way that is proportionate, consistent, and sustainable.}
\paragraph{Several elements of a coherent control framework are worth highlighting:}
\paragraph{\textbf{Policies and procedures} — the written documentation of expected behaviour and defined processes. Policies set the standards; procedures describe how those standards are implemented in practice. Both need to be current, accessible, and genuinely reflective of how the organisation operates — not aspirational documents that describe an idealised process bearing little resemblance to reality.}
\paragraph{\textbf{Delegated authorities and segregation of duties} — clear frameworks for who can authorise what, and the separation of functions that should not be combined in a single individual. Segregation of duties is one of the most effective administrative controls for preventing fraud and error — ensuring that no single individual has the ability to both initiate and approve a transaction, or to both create and validate a record.}
\paragraph{\textbf{Training and competence frameworks} — ensuring that the people responsible for operating controls have the knowledge, skills, and understanding to do so effectively. A policy that nobody understands is not a control. Training is not merely a box to be ticked — it is a fundamental component of control effectiveness, and its quality and relevance matter as much as its completion.}
\paragraph{\textbf{Management information and monitoring} — the data, reports, and oversight mechanisms that enable managers and the compliance function to maintain visibility of how controls are performing and whether the risk landscape is changing. We will explore monitoring in depth in Chapter~\ref{chap:Step 5 - Recording, Reviewing, and Monitoring}.}
\paragraph{\textbf{Escalation and reporting} — clear processes for ensuring that identified problems, control failures, and emerging risks are escalated promptly to the right level of management and, where necessary, to the board or to regulators.}

%---------------------------- SECTION ----------------------------
\section{Residual Risk — Accepting What Remains}
\paragraph{Once controls have been selected and designed, the organisation returns to the risk evaluation framework from Chapter~\ref{chap:Step 3 — Evaluating and Analysing Risks} to assess the \textbf{residual risk} — the level of risk that remains after the controls are in place and operating effectively.}
\paragraph{If the residual risk falls within the organisation's risk appetite, the control response is adequate and the risk can be accepted at the residual level, subject to ongoing monitoring.}
\paragraph{If the residual risk still exceeds the organisation's risk appetite after all proportionate controls have been implemented, the organisation faces a choice: accept the residual risk at a higher level than the stated appetite — with explicit sign-off from an appropriate level of authority — or revisit the fundamental question of whether to engage in the activity that creates the risk at all.}
\paragraph{This is the point at which risk appetite — discussed in Chapter~\ref{chap:Step 3 — Evaluating and Analysing Risks} — becomes operationally real. If the organisation's stated appetite for a particular risk type is genuinely zero, then an activity that cannot be controlled to zero residual risk should not be undertaken. If in practice the organisation is willing to accept some residual risk in that category, then the appetite statement needs to reflect reality — because a gap between stated appetite and accepted residual risk is a governance weakness that regulators and auditors will notice and challenge.}

%---------------------------- SECTION ----------------------------
\section{Documenting Control Decisions — The Evidential Record}
\paragraph{For compliance professionals, the documentation of control decisions is not merely an administrative convenience — it is a core component of the evidential record that demonstrates regulatory compliance. Regulators across virtually every sector expect to see not just that controls exist, but that they were deliberately chosen, that their selection was informed by a genuine risk assessment, and that the reasoning behind control decisions is recorded and defensible.}
\paragraph{Good documentation of control decisions includes:}
\begin{itemize}
    \bitem{A clear description of each control}{specific enough that an independent observer could verify whether it is being operated as described.}
    \bitem{The rationale for selecting each control}{why this control was chosen, and why it is considered proportionate to the risk}
    \bitem{The owner of each control}{the named individual accountable for its implementation and effectiveness.}
    \bitem{The implementation status}{is the control currently in place, being implemented, or planned?}
    \bitem{Any accepted residual risk}{where risk acceptance decisions have been made, the rationale and the level of authority at which the decision was made.}
    \bitem{The review date}{when the control and its effectiveness will next be formally assessed.}
\end{itemize}
\paragraph{This documentation feeds into and enriches the risk register developed in Chapter~\ref{chap:Step 3 — Evaluating and Analysing Risks} — adding the control dimension to each risk entry and creating a complete record of the organisation's risk management decisions.}

%---------------------------- SECTION ----------------------------
\section{Common Pitfalls in Control Design and Implementation}
\paragraph{As with the earlier steps, it is worth being explicit about the characteristic failure modes in this step.}

\subsection{Paper Controls}
\paragraph{Perhaps the most pervasive problem in compliance-driven risk management is the prevalence of \textbf{paper controls} — controls that exist in documentation but are not genuinely implemented in practice. A policy that was drafted three years ago and has not been reviewed since. A training programme that staff complete annually by clicking through slides without reading them. A management review process that appears on an organisational chart but has never actually taken place.}
\paragraph{Paper controls are dangerous precisely because they create the appearance of risk management without the substance. They may satisfy a superficial documentary review — but they will not withstand genuine scrutiny, and they will not prevent harm. The test of a control is not whether it exists on paper but whether it works in practice.}

\subsection{Control Clutter}
\paragraph{The opposite problem — accumulating controls without regularly reviewing whether they are still necessary, still proportionate, and still working — creates a different but related set of issues. Organisations that have grown their control frameworks incrementally over many years, adding new controls in response to each new risk or regulatory requirement without retiring obsolete ones, often end up with complex, overlapping, and contradictory control frameworks that are expensive to maintain and difficult for people to navigate.}
\paragraph{Periodic rationalisation of the control framework — removing controls that are no longer necessary, consolidating overlapping ones, and simplifying where complexity has accumulated without purpose — is good housekeeping that often also improves effectiveness.}

\subsection{Over-Reliance on Administrative Controls}
\paragraph{As noted in the discussion of the Hierarchy of Controls, administrative controls — policies, procedures, and training — are the most commonly deployed category of control but are inherently less reliable than higher-order controls. Organisations that rely almost exclusively on administrative controls, particularly for high-significance risks, are likely to find that those controls fail at precisely the moments when they are most needed — under pressure, when experienced staff are absent, or when the incentives to circumvent them are strongest.}
\paragraph{Where higher-order controls — elimination, substitution, or engineering controls — are available and practicable, they should be preferred.}

\subsection{Designing Controls Without Involving Those Who Will Operate Them}
\paragraph{Controls designed in isolation by a compliance or risk function, without genuine involvement of the operational teams who will be responsible for implementing and operating them, frequently suffer from the same basic problem: they do not reflect the operational reality of the processes they are designed to govern.}
\paragraph{The best-designed control in theory may be unworkable in practice because it conflicts with operational constraints, creates unacceptable friction, or relies on information or capabilities that operational teams do not have. Involving operational staff in control design — not just as recipients of the output but as genuine contributors to the process — produces better controls and significantly increases the likelihood that they will be properly implemented and maintained.}

\subsection{Treating Implementation as Completion}
\paragraph{Selecting and designing a control is the beginning of the work, not the end. Implementation takes time, requires active management, and frequently encounters obstacles that were not anticipated at the design stage. And once implemented, controls need to be monitored, tested, and maintained — because the environments in which they operate change, people change, and systems change.}
\paragraph{The most common single point of failure in the control lifecycle is the assumption that a control, once implemented, can be left to operate indefinitely without attention. It cannot. The monitoring and review step of Chapter~\ref{chap:Step 5 - Recording, Reviewing, and Monitoring} is the mechanism by which the organisation maintains assurance that its controls continue to work — and it is as important as the initial selection and implementation.}

%---------------------------- SECTION ----------------------------
\newpage
\section{Chapter Summary}
In this chapter we learned about the following key points:
\begin{table}[H]
  \centering
  \renewcommand{\arraystretch}{1.5}
  \begin{tabularx}{\textwidth}{ | >{\raggedright\arraybackslash}p{0.25\textwidth} 
                                | >{\raggedright\arraybackslash}X | }
    \hline
    \textbf{Response Option} & \textbf{When Most Appropriate} \\
    \hline
    \textbf{Avoid} & Risk is genuinely unacceptable and no other response can bring it within appetite\\
    \hline
    \textbf{Reduce} & Risk can be brought within appetite through proportionate prevention or mitigation controls\\
    \hline
    \textbf{Transfer} & Financial consequences can be shifted — but remember transfer does not remove regulatory or reputational risk\\
    \hline
    \textbf{Accept} & Risk falls within appetite — but acceptance must be conscious, documented, and authorised\\
    \hline
  \end{tabularx}
  \caption{Response Options}
\end{table}

\begin{table}[H]
  \centering
  \renewcommand{\arraystretch}{1.5}
  \begin{tabularx}{\textwidth}{ | >{\raggedright\arraybackslash}p{0.15\textwidth} 
                                | >{\raggedright\arraybackslash}p{0.35\textwidth} 
                                | >{\raggedright\arraybackslash}X | }
    \hline
    \textbf{Level} & \textbf{Type} & \textbf{Example in Compliance Context}\\
    \hline
    \textbf{1} & Elimination & Ceasing an activity that creates unacceptable risk\\
    \hline
    \textbf{2} & Substitution & Replacing a manual process with an automated one\\
    \hline
    \textbf{3} & System-enforced controls, mandatory approval workflows\\
    \hline
    \textbf{4} & Policies, procedures, training, management oversight\\
    \hline
    \textbf{5} & Detective & Monitoring, exception reporting, audit programmes\\
    \hline
    \textbf{6} & Protective & Business continuity plans, insurance, crisis response\\
    \hline
  \end{tabularx}
  \caption{Control Types}
\end{table}

\paragraph{Key takeaways:}
\begin{itemize}
  \item{The four fundamental risk response options — avoid, reduce, transfer, accept — are not mutually exclusive; most effective responses combine elements of several.}
  \item{The Hierarchy of Controls provides a framework for preferring more reliable control types where practicable — higher-order controls are inherently more robust than administrative ones.}
  \item{Proportionality is a key principle in control selection — the effort and cost of a control should reflect the significance of the risk it manages.}
  \item{Controls must be specific, owned, tested, and resilient to be genuinely effective — documented controls and working controls are not the same thing.}
  \item{Risk acceptance must be conscious, documented, and appropriately authorised — default acceptance is a gap, not a decision.}
  \item{The documentation of control decisions is a core component of the evidential record of regulatory compliance.}
  \item{Implementation is the beginning of the control lifecycle, not the end — monitoring and maintenance are essential and continuous.}
\end{itemize}

\barquote{In Chapter~\ref{chap:Step 5 - Recording, Reviewing, and Monitoring}, we reach the final step of the core process — exploring what good documentation looks like, how to build and maintain a risk register that is genuinely useful, how to set effective review cycles, and how ongoing monitoring keeps the risk assessment alive between formal reviews.}